% LaTeX Beamer 模板
% 根据具体项目修改标题、作者等信息
\documentclass[10pt,aspectratio=169]{beamer}

% 中文支持
\usepackage{ctex}

% 常用宏包
\usepackage{graphicx}
\usepackage{booktabs}
\usepackage{amsmath}
\usepackage{amssymb}
\usepackage{physics}
\usepackage{hyperref}
\usepackage{xcolor}
\usepackage{multirow}
\usepackage{algorithm}
\usepackage{algorithmic}
\usefonttheme[onlymath]{serif}

% 主题设置
\usetheme{Madrid}
\usecolortheme{default}

% 页边距设置
\setbeamersize{text margin left=1.2cm,text margin right=1.2cm}

% 图片路径（根据实际结构调整）
\graphicspath{{../images/}}

% 全局行间距
\renewcommand{\baselinestretch}{1.15}

% 引用设置
\setbeamertemplate{caption}[numbered]

% 每节开始时显示目录
\AtBeginSection[]
{
  \begin{frame}
    \frametitle{内容提要}
    \tableofcontents[currentsection]
  \end{frame}
}

% 标题页信息
\title[短标题]{完整标题}
\subtitle{副标题或会议名称}
\author[作者名]{汇报人：姓名}
\institute[机构简称]{
  机构全名或部门名称
}
\date{日期或事件}

\begin{document}

% 标题页
\frame{\titlepage}

% 目录页
\begin{frame}
  \frametitle{内容提要}
  \tableofcontents
\end{frame}

%=============================================================================
% 第一部分
%=============================================================================
\input{sections/part1_title.tex}

%=============================================================================
% 第二部分
%=============================================================================
\input{sections/part2_title.tex}

%=============================================================================
% 第N部分
%=============================================================================
\input{sections/partN_title.tex}

% 感谢页
\begin{frame}
  \centering
  \Huge{感谢聆听！}

  \vspace{1cm}
  \Large{欢迎提问与讨论}
\end{frame}

\end{document}
